\section{Introduction and relation to earlier work}\label{sec:intro}

There are two distinct uses of Loewner theory in a positivity construction.
One may try to identify a conformal-map observable directly with an
arithmetic transfer function. Alternatively, one may use the map to
generate a curve, evaluate a field-dependent observable on that curve,
and only then seek a transfer law. This paper develops the second option.
The additional field and endpoint data are essential: the conformal map
alone does not determine Wilson transport or its variation.

The proposed application is the Riemann Hypothesis (RH). In the
normalization used here, RH is equivalent to positivity of an explicit
quadratic form on every compact interval. Its archimedean part contains
the kernel
\begin{equation}\label{eq:intro-tower}
 \nGamma(u)=\frac{e^{-u/2}}{1-e^{-2u}}
 =\sum_{n=0}^{\infty}e^{-(2n+1/2)u},\qquad u>0.
\end{equation}
A free three-dimensional defect scalar, after a logarithmic radial
change of variables and even angular smearing, supplies this sequence
of exponents. That observation motivates the use of defect matter.
It does not supply the local subtraction, the pole terms, or the
prime-power terms of the Weil form. Section~\ref{sec:arithmetic} states
that target explicitly, and Section~\ref{sec:free} explains exactly what
the free construction provides.

\subsection{The relevant literature}

The scalar-coupled Wilson operator in four-dimensional
$\mathcal N=4$ supersymmetric Yang--Mills theory is part of the
Wilson-loop framework introduced in the AdS/CFT setting by
Maldacena~\cite{Maldacena}. Zarembo~\cite{Zarembo} constructed families
whose internal scalar direction follows the spacetime tangent through
a constant isometry. The resulting linear Clifford projectors allow
arbitrary contour shapes while retaining common supersymmetry.
Our tangent prescription is an application of that construction. The
additional questions here are its compatibility with the chiral defect
projector, physical endpoint matter, and reflection.

DeWolfe, Freedman and Ooguri~\cite{DFO} described a four-dimensional
$\mathcal N=4$ bulk theory coupled to a three-dimensional fundamental
hypermultiplet at a D3--D5 intersection. The defect preserves an
$SU(2)_H\times SU(2)_V$ R symmetry. Gaiotto and Witten~\cite{GW} give
the associated brane supersymmetry projectors in their study of
supersymmetric boundary conditions. We use the unfluxed intersection
projector, not their additional conditions for D3-branes terminating on
a D5-brane. Baker~\cite{Baker} constructed supersymmetric open Wilson
lines in this defect theory, including a semicircle obtained from a ray
by a conformal transformation, and analyzed their scalar endpoints.
These examples are controls for our endpoint test with a constant
Poincare charge; that test does not exclude Baker's semicircle.

Loewner's equation describes the conformal maps of growing planar hulls.
We use the chordal normalization in Lawler's lectures~\cite{Lawler}.
The relation between stochastic Loewner evolution and conformal-field-
theory martingales developed by Bauer and Bernard~\cite{BauerBernard}
is relevant motivation. A map-level stochastic differential equation,
however, does not by itself construct a renormalized Wilson observable
on a rough trace. We work with smooth curves and regular pullbacks;
the one stochastic formula below is explicitly limited to the inverse map.

Osterwalder--Schrader reflection positivity~\cite{OS} explains how
Euclidean reflected pairings can define Hilbert-space norms. For
scalar-coupled Wilson transport, reflection must account for both
complex conjugation and path ordering. Dorn and Pershin~\cite{DP}
analyzed this distinction and separated ordinary reflection inequalities
from a conjectural inequality involving an internal rotation in a
special large-rectangle limit. Our internal-rotation counterexample
concerns the full defect scalar doublet and does not address that
restricted conjecture.

On the arithmetic side, Weil's positivity criterion provides the
logical target. Suzuki's recent treatment~\cite{Suzuki} places the
localized Weil form, the screw function and associated operator
constructions in a common framework, with references to the earlier
work of Weil, Yoshida, Bombieri, and Connes and collaborators.
We use the compact-support criterion, not a conjectural limiting
spectral construction. The present field-theoretic calculations establish
no new equivalence with RH.

\subsection{What survives from the parent investigations}

The preceding Wilson-lines, Loewner, and fractional-dimension
investigations~\cite{ParentWilson,ParentLoewner,ParentFractional} are
working research documents in the same program. Their relevant
identities are reproduced here, so they are not prerequisites for
reading this paper. In particular, the shifted transfer has an explicit
positive Beta-kernel/prime-comb part and an essential rational correction.
The earlier direct Loewner realization tests did not identify a bare
map observable producing this entire arithmetic object. The
fractional-dimension calculation explains the gamma ratio by a
dimensionally continued integral, but a natural continued lattice count
has negative coefficients. These are restrictions on particular
realizations, not an exclusion of field-dependent observables on
Loewner-generated curves.

The combined construction studied here therefore has four separate
ingredients: a contour law; a gauge and scalar connection on that
contour; endpoint and reference states with a specified adjoint; and
an eventual arithmetic identity. We establish parts of the first three.
We do not identify the fourth.

\subsection{Results and scope}

The argument proceeds through the following results.
\begin{enumerate}
 \item Inversion sends the inherited common-origin semicircles to
 vertical slits of constant-driver chordal evolution. Under the fixed
 scalar and special bulk charge specified in
 Section~\ref{sec:geometry}, those curves are rigid.
 \item The exact smooth variation of an open Wilson transport contains
 field-strength, scalar-gradient, arclength, and endpoint terms. An
 abelian counterexample rules out a universal geometry-only scalar
 closure of the transport variation (Section~\ref{sec:variation}).
 \item A real tangent isometry compatible with chirality and the D3--D5
 projector admits a two-dimensional complex charge space in a normal
 plane, and a one-dimensional space across all defect directions.
 Its orientation matters (Section~\ref{sec:projectors}).
 \item Ordinary scalar endpoints annihilated by the same constant charge
 form a zero-contraction pair. The full bilinear has nonzero residual
 R charge and zero expectation in an invariant vacuum. Physical adjoints
 evade the zero contraction but change the supersymmetry question
 (Section~\ref{sec:endpoints}).
 \item Ordinary time reflection gives a different scalar map on the
 reflected branch. The resulting free kernel is nonzero and positive.
 Positivity of the interacting network is stated conditionally, with
 its color and regulator requirements (Section~\ref{sec:reflection}).
 \item A family with a smooth reference join has a finite direct
 Gaussian bulk response with leading coefficient
 $3/2-2\log2>0$. Endpoint, defect and renormalization contributions
 have not been computed (Section~\ref{sec:response}).
\end{enumerate}

All supersymmetry ranks are complex ranks in a specified Euclidean
Clifford module. All transport variation theorems are statements about
smooth finite matrix fields. A regulated quantum use of those identities
requires additional analytic and renormalization input. The finite
diagnostics in Appendix~\ref{app:diagnostics} check signs, ranks and
integrals; they neither establish a continuum quantum field theory nor
prove an inequality for the Weil form.
